\chapter{电磁加速与温度灵敏度：怎样提高阻抗精度}
\label{ch:tangent}

\section{为什么在神经网络之外保留物理校正}
第~\ref{ch:neural} 章的记忆网络可以快速估计动态海水复阻抗。面对训练分布以外的线圈姿态、封装尺寸或高功率温度变化，海水电磁反作用对局部温度的位置可能很敏感。若只有网络预测，误差可能明显增大。

我们已经有原 BFZI 这一成熟的电磁求解器，因此采取一条计算量较小的物理途径：在每个新几何的初始温度附近，提前算好阻抗对若干代表性海水温度形状的变化率。在线时只需将当前温度场投影到这些形状上，再用预先存好的变化率估计海水阻抗变化。

这种变化率称为\textbf{温度灵敏度}，也称温度方向导数或温度切线。它是完全可以用物理计算解释的量。

\section{先用八个空间温度模式描述海水温变}
海水内部可能有成千上万个温度节点。直接对每一个节点分别求阻抗导数，离线计算量很大。我们从热 ROM 的空间模态中提取八个海水温度模式
\[
B_1(\mathbf x),\dots,B_8(\mathbf x),
\]
将当前温升 $\theta(\mathbf x,t)=T(\mathbf x,t)-T_0$ 写成
\[
\theta(\mathbf x,t)\approx \sum_{j=1}^8\beta_j(t)B_j(\mathbf x)
+\theta_\perp(\mathbf x,t).
\]
其中 $\theta_\perp$ 是这八个模式未能描述的剩余空间温度变化。模式在海水有限体积上按体积正交，系数为
\begin{equation}
\beta_j(t)=\sum_i V_i B_j(\mathbf x_i)\theta_i(t).
\end{equation}

八个温度模式可以理解为八种有代表性的加热形状。例如一部分模态对应整体海水升温，另一部分模态描述靠近某只封装的温度梯度。模型实际从物理热响应自动提取它们，而不是手工指定热点形状。

\section{一个模式改变一点温度，阻抗如何变化}
对第 $j$ 个模式，先在冷态温度 $T_0$ 的海水场上增加少量扰动 $\delta B_j$，用原 BFZI 计算新的复海水阻抗。用
\begin{equation}
\mathbf D_j=
\left.\frac{\partial\zsea}{\partial\beta_j}\right|_{T_0}
\end{equation}
表示沿此模式的阻抗变化率。$\mathbf D_j$ 是一个 $2\times2$ 的完整复矩阵，其中每个独立元素都包含实部与虚部。

在线预测为
\begin{equation}
\label{eq:tangent}
\Delta\zsea^{\rm phys}(t)
=\sum_{j=1}^8\beta_j(t)\mathbf D_j.
\end{equation}
每个时间步只包含八个系数和八个小矩阵的乘加，不需要再解海水 Green/resolvent 方程。

\section{八次和十六次电磁方向求解}
早期实现使用双侧差分：
\[
\mathbf D_j^{\rm central}\approx
\frac{\zsea(T_0+\delta B_j)-\zsea(T_0-\delta B_j)}
{2\delta}.
\]
每个模式需要两个原 BFZI 电磁求解，共 $8\times2=16$ 次。

后来加入单侧差分：
\[
\mathbf D_j^{\rm forward}\approx
\frac{\zsea(T_0+\delta B_j)-\zsea(T_0)}{\delta}.
\]
由于冷态 $\zsea(T_0)$ 已存在冷态锚中，每个模式只需再求一次，八个模式共 8 次。历史实际实验在 P03/P07 的相同温度模式上确认：单侧差分使用约 0.2\,K 的方向温度扰动，完整动态复阻抗精度与双侧方案十分接近，方向求解计算量约下降一半。

\begin{table}[htbp]\centering\small
\caption{P07 真实 150\,V、15\,s 独立时程中不同切线方案的动态完整复阻抗误差}
\begin{tabularx}{\textwidth}{Xcc}
\toprule
计算方法 & ID16 & ID18\\
\midrule
双侧十六次原 BFZI 方向求解 & 0.001902\% & 0.002579\%\\
单侧八次原 BFZI 方向求解 & 0.001924\% & 0.002606\%\\
跨几何零新增方向、利用已有时程回归与迁移 & 0.211073\% & 0.253541\%\\
\bottomrule
\end{tabularx}
\end{table}

\section{为什么完整复阻抗误差可以很小}
表中的 $0.002\%$ 量级是\textbf{完整复阻抗的动态增量}误差。其中铜温变部分由原 BFZI 的精确铜阻抗公式计算，物理切线主要校正海水温变部分。因此完整复阻抗误差与单独海水 $\Delta\zsea$ 的相对误差是不同指标。

在 P03/P07 的早期实验中，八模态海水增量自身的相对误差约为 0.5\%--1.3\%，而完整动态复阻抗由于还包含精确铜阻抗变化，其相对误差可达到约 0.002\%。科研报告应同时列出两种指标，才能判断海水校正本身的质量。

\section{神经预测与物理切线怎样组合}
给定网络预测 $\Delta\zsea^{\rm NN}$ 与式~\eqref{eq:tangent} 的物理预测，定义
\begin{equation}
\label{eq:blend}
\Delta\zsea=
(1-\lambda)\Delta\zsea^{\rm NN}
+\lambda\Delta\zsea^{\rm phys},\qquad 0\le\lambda\le1.
\end{equation}
$\lambda=0$ 使用纯神经海水增量；$\lambda=1$ 使用物理切线预测；中间值作加权组合。

在 \code{GeomHeat.rollout} 中，算法先用神经阻抗作一个热态预测，再在\emph{预测的新三维温度}上评价物理切线并作正式热步修正。这样可以避免只用上一热步旧温度来计算新阻抗的时间滞后。

当前历史测试中最高精度模式使用 $\lambda=1$。神经网络仍为热态预测提供先验，温度相关海水阻抗的最终读出由物理切线提供。

\section{电磁 ROM 与电磁切线有什么区别}
仓库早期还实现过\textbf{电磁 ROM}：对多个原 BFZI 单位端口场快照做低秩投影，将大规模电磁线性方程降为小规模方程，每个热步在当前电导率下更新小矩阵求解。这条路线保持原电磁方程结构，早期 M4 留出精度很好，但小样本上的整体加速相对有限。

当前高精度动态路线则\textbf{预计算温度切线}，在线以八个温度系数直接得到海水阻抗变化。二者都利用低维物理空间，计算对象不同：电磁 ROM 降低每次求解电磁方程的成本，温度切线减少每步对电磁方程的重新求解。

\section{跨几何共享：怎样减少新几何的准备成本}
对一个从未出现过的几何，重新建立冷态场、热算子、八个温度方向导数仍然需要时间。因此代码还实现了跨几何共享 ROM 路线。

在若干已知几何上保存物理热 ROM、海水温度切线和冷态场修正。新几何只执行一次原 BFZI 冷态求解，利用与它相近的供体几何提供热基和切线，并对冷态端口进行修正。实验中的 P03/P07/G3 整几何留出使用该方法，动态复阻抗误差约为 0.83\%--1.31\%，新几何离线成本下降约 58\%--63\%。

此外还有跨几何\textbf{神经切线迁移}：先对齐不同几何下热模态的次序、符号和旋转，再学习八模态复阻抗导数之间的变化。已实现的算法使用体积加权重叠和正交 Procrustes 配准；零新增方向电磁求解的精度与八次物理校准路径之间已有实际对照。

\section{如何选择}
当研究目标是\textbf{最高动态复阻抗精度}时，优先使用新几何的八模态物理校准与原三维热 ROM。当目标是\textbf{大量相近新几何的低成本部署}时，采用共享 ROM 或神经切线迁移，记录精度与离线成本的变化。两条路径都保留 BFZI 冷态端口与三维物理场信息，区别在于需要为目标几何重新计算多少温度相关物理量。

相关代码为 \code{geom\_tangent.py}、\code{geom\_tangent\_transfer.py}、\code{geom\_shared\_rom.py} 和 \code{v148\_rom.py}。
